\section{Introduction}
\label{sec:intro}

Coding agents are now routinely asked to modify real software. Most evaluations
of that ability measure whether an agent can make a failing test pass
\cite{jimenez2024swebench} --- a task with a clear oracle, where success is
checkable by running the suite. Performance optimization has no such oracle. An
agent asked to make a program faster must first decide \emph{what} to make
faster, and a plausible-looking change that helps nothing is indistinguishable
from a real improvement unless someone measures carefully.

This paper asks a narrow version of that question. Given a mature C codebase, no
hint about where time is spent, and an instruction only to make the program
faster without changing its output, can a coding agent find a genuine bottleneck
and fix it?

We use the \texttt{dot} layout engine from GraphViz~\cite{ellson2001graphviz}.
It is real, widely deployed software rather than a benchmark written for the
occasion; its layered pipeline
\cite{sugiyama1981methods,gansner1993technique} has distinct phases whose costs
can be measured separately, so a speedup can be attributed rather than merely
observed; and its text output is deterministic, which makes ``the output did not
change'' checkable rather than a judgment call.

\paragraph{What we did.}
Two agents --- Claude Code and Codex --- each received a pristine checkout of
GraphViz pinned to commit \texttt{23b06521} and an identical prompt naming no
target function, file, or phase. The prompt explicitly forbade the shortcuts
that would fake a win: caching layout results, special-casing inputs, and
reducing the iteration counts that govern layout quality. Neither agent had
access to the evaluation harness or the graphs it uses. We rebuilt each agent's
result ourselves and measured it against an untouched baseline, on three graph shapes with the node count held fixed.

\paragraph{What we found.}
Both agents independently identified crossing minimization as the target, and
both produced byte-identical layout output.

They separate on breadth. Claude improved three phases --- ranking, crossing
minimization, coordinate assignment --- measuring $1.33\times$ to $1.84\times$
end-to-end. Codex improved crossing minimization only, measuring $1.29\times$
down to $0.99\times$: on one shape its optimization is real but the phase is
$11.5\%$ of the workload, so nothing survives to the end-to-end result.

Which phase dominates is a property of the graph's \emph{shape}, not its size.
At identical node count, crossing minimization ranges from $11.5\%$ to $64\%$ of
layout time, setting a different ceiling on each workload --- and on the sparsest
one Claude exceeds the ceiling available to any crossing-minimization
optimization at all.

\paragraph{Contributions.}
\begin{itemize}
\item A controlled protocol for evaluating agent-driven optimization: balanced
      command ordering, which also yields an empirical noise floor from the
      baseline's own spread across orderings, and a held-out evaluation
      harness.
\item Phase-resolved measurements of two agents' work across three graph
      shapes, with per-phase attribution and the shape-dependent ceiling
      that bounds it.
\item A complete artifact: pinned source, both agents' diffs and commit
      histories, raw timing data, and the figure-generation scripts.
\end{itemize}

Section~\ref{sec:related} covers background. Section~\ref{sec:method} describes
the protocol. Section~\ref{sec:results} reports
measurements. Sections~\ref{sec:discussion}--\ref{sec:limitations} interpret
them and state what we did not verify.
